\section{(H) 函数的反函数}

若 \(f\) 是一个 (H) 函数，则 \(f\) 在某个区间 \([x_0, +\infty[\) 上单调且连续，故限制的反函数 \(\varphi\) （对 \(f\) 在该区间上的限制）在点 \(a = \lim_{x \to +\infty} f(x)\) 的一个邻域上单调且连续；但若 \(a\) 等于 \(+\infty\) （相应地， \(-\infty\) 、有限），则可以证明 \(\varphi(y)\) （相应地， \(\varphi(-y)\) 、 \(\varphi(a + \frac{1}{y})\) 或 \(\varphi(a - \frac{1}{y})\) ）在 \(+\infty\) 的某个邻域上一般不等于一个 (H) 函数。尽管如此，我们将看到，在某些重要情形可以得到与 \(\varphi(y)\) （相应地， \(\varphi(-y)\) 、 \(\varphi(a + \frac{1}{y})\) 、 \(\varphi(a - \frac{1}{y})\) ）等价的 (H) 函数，有时甚至可得到这个函数相对于 V, p.~254 定义的尺度 \(\mathcal{E}\) 的渐近展开。

我们将用到下列命题：

命题 4. 设 \(p\) 与 \(q\) 是两个 (H) 函数（在某个区间 \([x_0, +\infty[\) 上严格正）。

\(1^{\circ}\) 若 \(q \ll p / p'\) ，则 \(p(x + q(x)) \sim p(x)\) 。

\(2^{\circ}\) 若 \(q \prec p/p'\) 且 \(q(x) \prec x\) ，则 \(p(x - q(x)) \sim p(x)\) 。

若 \(p \sim k\) （常数 \(\neq 0\) ），这个命题的两部分都是显然的；故可设 \(p(x) \ll 1\) （否则把论证应用于 \(1/p\) ）。由此推出 \(p'(x) \ll 1\) 。

\(1^{\circ}\) 可以写 \(p(x + q(x)) = p(x) + q(x)p'(x + \theta q(x))\) ，其中 \(0 \leqslant \theta \leqslant 1\) (I, p.~14, corollary)。由于 \(\left|p'(x)\right|\) 当 \(x\) 趋于 \(+\infty\) 时趋于 0，且在 \(+\infty\) 的某个邻域上等于一个 (H) 函数，故它在某个区间 \([x_1, +\infty[\) 上递减，于是对 \(x \geqslant x_1\) 有 \(\left|p'(x + \theta q(x))\right| \leqslant \left|p'(x)\right|\) ；由于 \(qp' \ll p\) ，得 \(p(x + q(x)) \sim p(x)\) 。

\(2^{\circ}\) 条件 \(q(x) \ll x\) 保证 \(x - q(x)\) 趋于 \(+\infty\) （随 \(x\) ）。仍有 \(p(x - q(x)) = p(x) - q(x)p'(x - \theta p(x))\) ，其中 \(0 \leqslant \theta \leqslant 1\) 。与证明第一部分同样的论证表明，对充分大的 \(x\) ， \(\left|p'(x - \theta q(x))\right| \leqslant \left|p'(x - q(x))\right|\) 。于是归结为证明 \(q(x)\frac{p'(x - q(x))}{p(x - q(x))}\) 当 \(x\) 趋于 \(+\infty\) 时趋于 0。若 \(p'/p \gg 1\) ，命题成立，因为这时 \(\left|p'/p\right|\) 是一个 (H) 函数，且对充分大的 \(x\) 递增；于是 \(q(x)\frac{\left|p'(x - q(x))\right|}{\left|p(x - q(x))\right|} \leqslant q(x)\frac{\left|p'(x)\right|}{\left|p(x)\right|}\) ，而由假设 \(qp' \ll p\) 。若 \(p'/p \sim k\) （ \(k\) 为常数 \(\neq 0\) ），命题也成立，因为这时

\[
\frac {p ^ {\prime} (x - q (x))}{p (x - q (x))} \sim \frac {p ^ {\prime} (x)}{p (x)}
\]

因为 \(x - q(x)\) 趋于 \(+\infty\) 。只剩考察 \(p' / p \ll 1\) 的情形。首先设 \(p(x)\) 相对于 x 的阶为有限，于是 (V, p.~232, prop. 7) \(p'(x) / p(x) \ll 1 / x\) 。这时有

\[
\frac {p ^ {\prime} (x - q (x))}{p (x - q (x))} = \frac {1}{x - q (x)} O _ {1} (1),
\]

于是

\[
q (x) \frac {p ^ {\prime} (x - q (x))}{p (x - q (x))} = \frac {q (x)}{x} \left(1 - \frac {q (x)}{x}\right) ^ {- 1} O _ {1} (1) = \frac {q (x)}{x} O _ {2} (1),
\]

于是在这一情形可以看出，仅在不等式 \(q(x) \ll x\) 的假设下命题就成立。最后我们考察 \(1/x \ll p'(x)/p(x) \ll 1\) 的情形；这时函数 \(r = p'/p\) 相对于 x 的阶为有限；由于按前面的注，V, p.255 的命题 4 可应用于这样的函数，有 \(p'(x-q(x))/p(x-q(x)) \sim p'(x)/p(x)\) ，而假设 \(qp' \ll p\) 使我们能完成证明。

注. 加在 \(q(x)\) 上的条件不能改进，如下例所示：

\begin{enumerate}
\def\labelenumi{\alph{enumi})}
\tightlist
\item
\end{enumerate}

\[
p (x) = e ^ {x}, \qquad q (x) = 1 = \frac {p (x)}{p ^ {\prime} (x)}, \qquad p (x + q (x)) = e p (x)
\]

\begin{enumerate}
\def\labelenumi{\alph{enumi})}
\setcounter{enumi}{1}
\tightlist
\item
\end{enumerate}

\[
\begin{array}{c} p (x) = \log x, \qquad q (x) = x - \log x \ll \frac {p (x)}{p ^ {\prime} (x)} = x \log x, \\ p (x - q (x)) = \log \log x \ll p (x). \end{array}
\]

我们先研究一类特殊 (H) 函数的反函数：

命题 5. 设 g 是一个不等价于常数 \(\neq 0\) 的 (H) 函数，且满足 \(g(x) \prec x\) ，并设 \(u(x)\) 是 \(x - g(x)\) 的反函数（在 \(+\infty\) 的某个邻域上有定义）。设 \((u_n)\) 是由条件 \(u_0(x) = x\) ， \(u_n(x) = x + g(u_{n-1}(x))\) （对 \(n \geqslant 1\) ）对 n 递归定义的函数序列；则 \(u_n \gg 1\) 且

\[
u (x) - u _ {n} (x) \sim g (x) \left(g ^ {\prime} (x)\right) ^ {n}.\tag{11}
\]

令 \(y = u(x)\) ， \(y_{n} = u_{n}(x)\) ；则 \(x = y - g(y)\) ， \(y_{0} = x\) 且 \(y_{n} = x + g(y_{n+1})\) 。首先推出 \(x/y = 1 - \frac{g(y)}{y}\) ；由于 y 随 x 趋于 \(+\infty\) ，假设 \(g(x) \ll x\) 表明 \(y = u(x) \sim x = y_{0}\) ；此外，

\[
y - x = g (y) = g (x) + (y - x) g ^ {\prime} (z)
\]

其中 z 属于以 x, y 为端点的区间；当 x 趋于 \(+\infty\) 时 z 亦然，且由于 \(g(x) \ll x\) ， \(g' \ll 1\) ，故 \(g'(z)\) 趋于 0，从而

\[
y - x = g (x) + o (y - x)
\]

由此得

\[
u (x) - x \sim g (x).\tag{12}
\]

其次我们用对 \(n\) 的递归来证明：当 \(x\) 趋于 \(+\infty\) 时有 \(u_{n} \gg 1\) ，且

\[
u (x) - u _ {n} (x) \prec u (x) - u _ {n - 1} (x).\tag{13}
\]

事实上， \(y - y_{n} = g(y) - g(y_{n-1}) = (y - y_{n-1})g'(z_{n-1})\) ，其中 \(z_{n-1}\) 属于以 y 与 \(y_{n-1}\) 为端点的区间；按归纳假设 \(z_{n-1}\) 随 x 趋于 \(+\infty\) ，故 \(g'(z_{n-1})\) 趋于 0，这证明了 (13)。由这个关系式与 (12) 推出 \(u(x) - u_{n}(x) \ll u(x) - x \sim g(x) \ll x \sim u(x)\) ，于是 \(u_{n}(x) \sim u(x)\) ，从而 \(u_{n} \gg 1\) 。最后，关系式 \(u(x) - u_{n}(x) \ll u(x) - x\) 也可写成 \((u(x) - x) - (u_{n}(x) - x) \ll u(x) - x\) ，由此得

\[
u _ {n} (x) - x \sim u (x) - x \sim g (x).\tag{14}
\]

为证 (11)，首先注意：若 \(t(x)\) 是满足 \(t(x) - x \sim g(x)\) 的函数，则有 \(g'(t(x)) \sim g'(x)\) 。事实上，对任意 \(\varepsilon > 0\) ，当 x 充分大时 \(g'\) 单调，故 \(g'(t(x))\) 介于 \(g'(x + (1 + \varepsilon)g(x))\) 与 \(g'(x + (1 - \varepsilon)g(x))\) 之间。V, p.~255 的命题 4 表明 \(g'(t(x)) \sim g'(x)\) ，前提是已建立关系式 \(g \ll g'/g''\) 。而若 g 相对于 x 的阶为无穷，则有 (V, p.~251, prop. 2) \(g''/g' \sim g'/g\) ，且由于 \(g' \ll 1\) ，有 \(g \ll g/g' \sim g'/g''\) ；若 g 相对于 x 的阶为有限数 \(\mu\) ，则必有 \(\mu \leqslant 1\) ；若 \(\mu < 1\) ，则由于 g 不等价于常数 \(\neq 0\) ，公式 (3) 与 (4) (V, p.~251) 表明 \(g''/g' \sim k/x\) （k 为常数 \(\neq 0\) ），于是仍有 \(g \ll g'/g''\) ；最后，若 \(\mu = 1\) ，则 \(g'\) 相对于 x 的阶为 0，故 \(g''/g' \ll 1/x\) ，从而仍有 \(g \ll g'/g''\) 。

这样，由于 \(z_{n-1}\) 介于 \(y\) 与 \(y_n\) 之间，由 (14) 得 \(z_{n-1} - x \sim g(x)\) ，从而由上所述 \(g'(z_{n-1}) \sim g'(x)\) ；于是

\[
y - y _ {n} \sim (y - y _ {n - 1}) g ^ {\prime} (x),
\]

再对 \(n\) 作归纳便得 (11)。

注. 1) 若 \(g\) 的阶 \(< 1\) （相对于 \(x\) ），则函数 \(u(x) - u_n(x)\) 随 \(x\) 趋于 0，只要 \(n\) 充分大。事实上，在相反的情形将有 \(gg'^n \gg 1\) 对每个 \(n\) 成立，于是 \(g\) 相对于 \(1 / g'\) 的阶为无穷；换言之，将有 \(\log |g| \gg \log |g'|\) ，求导便得 \(g'/g \gg g''/g'\) 。但若 \(g\) 的阶为 \(\mu < 1\) （相对于 \(x\) ），则有 \(g'/g \sim g''/g'\) （当 \(\mu = -\infty\) ）， \(\frac{g'}{g} \sim \frac{\mu}{\mu - 1} \frac{g''}{g'}\) （当 \(\mu \neq 0\) ），最后 \(g'/g \ll g''/g'\) （当 \(\mu = 0\) ）(V, p.~251, n° 3)。

相反，若 \(g\) 相对于 \(x\) 的阶为 1，则可能 \(gg'^n \gg 1\) 对每个整数 \(n > 0\) 成立，如例子 \(g(x) = x / \log x\) 所示。

\begin{enumerate}
\def\labelenumi{\arabic{enumi})}
\setcounter{enumi}{1}
\tightlist
\item
  当 \(g(x)\) 是等价于常数 \(k \neq 0\) 的 (H) 函数时，有 \(g(x) = k + g_1(x)\) ，其中 \(g_1 \ll 1\) ；函数 \(u_1(x) = u(x) - k\) 是函数 \(x - g_1(x + k)\) 的反函数，于是归结为 V, p.~256 命题 5 所处理的情形。
\end{enumerate}

于是，要得到函数 u 的渐近展开，只需得到函数 \(u_{n}\) 的渐近展开：若 g 相对于所考虑的尺度有渐近展开，则（由 (H) 函数的定义）这就归结为 V, p.~223 至 227 中考察过的问题。

如下最一般的情形可归结为 V, p.~256 命题 5 所处理的情形：设函数 \(y = u(x)\) 满足关系式

\[
\varphi (x) = \psi (y) - g (y)\tag{15}
\]

其中 \(\varphi\) 是一个 (H) 函数， \(\psi\) 是满足 \(\psi \gg 1\) 、且 \(\theta\) （ \(\psi\) 的反函数）也是 (H) 函数的 (H) 函数，而 g 是满足 \(g \ll \psi\) 的 (H) 函数。现设 \(v(x)\) 是 \(x - g(\theta(x))\) 的反函数；则有 \(u = \theta \circ v \circ \varphi\) ，且 \(g(\theta(x)) \ll x\) ；若借助 V, p.~256 的命题 5 已知 v 的一个渐近展开，则可如 V, p.~223 至 227 中那样推出 u 的渐近展开。

例. 1) 试求反函数 \(v(x)\) （ \(x^{5} + x\) 的反函数）的渐近展开（对趋于 \(+\infty\) 的 x）；令 \(x^{5} = t\) ，归结为求反函数 \(u(t)\) （ \(t + t^{1/5}\) 的反函数）的展开（对趋于 \(+\infty\) 的 t），也就是把 V, p.~256 的命题 5 应用于 \(g(t) = -t^{1/5}\) 的情形。例如，计算 \(u_{2}(t)\) ；有

\[
u _ {2} (t) = t - \left(t - t ^ {1 / 5}\right) ^ {1 / 5} = t - t ^ {1 / 5} + \frac {1}{5} t ^ {- 3 / 5} + \frac {2}{2 5} t ^ {- 7 / 5} + o _ {1} \left(t ^ {- 7 / 5}\right).
\]

此外，由 (11) (V, p.~256) 有

\[
u (t) - u _ {2} (t) \sim - \frac {1}{2 5} t ^ {- 7 / 5}
\]

由此得

\[
u (t) = t - t ^ {1 / 5} + \frac {1}{5} t ^ {- 3 / 5} + \frac {1}{2 5} t ^ {- 7 / 5} + o _ {2} \left(t ^ {- 7 / 5}\right)
\]

于是推出展开

\[
v (x) = u (x) ^ {1 / 5} = x ^ {1 / 5} - \frac {1}{5} x ^ {- 3 / 5} - \frac {1}{2 5} x ^ {- 7 / 5} + o _ {3} \left(x ^ {- 7 / 5}\right)
\]

即为所求。

\begin{enumerate}
\def\labelenumi{\arabic{enumi})}
\setcounter{enumi}{1}
\tightlist
\item
  试求函数 \(x/\log x\) 的反函数的渐近展开；由恒等式 \(x = y/\log y\) （其中 \(y = v(x)\) ）推出 \(\log x = \log y - \log \log y\) ；令 \(z = \log y\) ， \(t = \log x\) ，有 \(t = z - \log z\) ，于是归结为展开反函数 \(u(t)\) （ \(t - \log t\) 的反函数）；例如，有
\end{enumerate}

\[
u _ {2} (t) = t + \log (t + \log t) = t + \log t + \frac {\log t}{t} - \frac {(\log t) ^ {2}}{2 t ^ {2}} + o _ {1} \left(\frac {\log t}{t ^ {2}}\right)
\]

此外，由 (11) (V, p.~256) 有

\[
u (t) - u _ {2} (t) \sim \frac {\log t}{t ^ {2}}
\]

由此得

\[
u (t) = t + \log t + \frac {\log t}{t} - \frac {(\log t) ^ {2}}{2 t ^ {2}} + \frac {\log t}{t ^ {2}} + o _ {2} \left(\frac {\log t}{t ^ {2}}\right)
\]

回到原来的问题，便得到渐近展开

\[
v (x) = x \log x + x \log \log x + x \frac {\log \log x}{\log x} + o \left(x \frac {\log \log x}{\log x}\right).
\]

\pandocbounded{\includegraphics[keepaspectratio]{images/part002_15bf1a92.jpg}}

注. 注意两个等价的 (H) 函数可以有不等价的反函数，如两个函数 \(\log x\) 与 \(1 + \log x\) 的例子所示。

\setcounter{exercisegroup}{0}
\refstepcounter{exercisegroup}
